\documentclass[12pt,a4paper]{article}
\usepackage[utf8]{inputenc}
\usepackage[spanish,es-noquoting]{babel}
\usepackage[T1]{fontenc}
\usepackage{lmodern}
\usepackage{geometry}
\geometry{top=2cm,bottom=2cm,left=2.5cm,right=2.5cm}
\usepackage{xcolor}
\usepackage{listings}
\usepackage{enumitem}
\usepackage{fancyhdr}
\usepackage{graphicx}
\graphicspath{{../../assets/}}
\lstset{language=C,basicstyle=\ttfamily\small,breaklines=true,frame=single,showstringspaces=false}
\pagestyle{fancy}
\setlength{\headheight}{15pt}
\fancyhf{}
\lhead{Monitorías Sistemas Operativos 2026-II}
\rhead{Capítulo 3: Hilos}
\cfoot{\thepage}
\title{\textbf{Taller Propuesto: Procesamiento Concurrente con Hilos}}
\author{Monitorías Sistemas Operativos}
\date{\today}
\begin{document}
\maketitle

\section*{Propósito}
Los ejercicios integran creación de hilos, paso de argumentos, memoria compartida y recolección de resultados con \texttt{pthread\_join}. Compile con \texttt{-pthread}.

\section*{Ejercicio 1}
Un archivo contiene primero $N$ y después $N$ mediciones enteras. El programa recibe $P$ por línea de comandos y divide el vector en $P$ bloques contiguos.
\begin{enumerate}
    \item Cree un hilo por bloque y páselo mediante una estructura que incluya inicio, fin y puntero al vector.
    \item Cada hilo debe calcular la suma y el máximo de su bloque sin imprimir desde la función de hilo.
    \item El \texttt{main} debe hacer \texttt{join} a todos los hilos, consolidar la suma y el máximo global e imprimir ambos resultados.
    \item Considere el caso en que $N$ no sea divisible entre $P$.
\end{enumerate}

\section*{Ejercicio 2}
Reciba un archivo de texto y un número $P$. El programa debe cargar cada palabra en memoria dinámica y distribuir rangos de palabras entre $P$ hilos.
\begin{enumerate}
    \item Cada hilo cuenta cuántas palabras de su rango comienzan con vocal y cuántas superan seis caracteres.
    \item Devuelva ambos conteos mediante la estructura de argumentos, no mediante variables globales.
    \item Después de todos los \texttt{join}, imprima los conteos totales y libere cada cadena y el vector de punteros.
\end{enumerate}

\section*{Ejercicio 3}
Se dispone de una estructura compartida con dos campos: \texttt{int aprobados} y \texttt{int rechazados}. Cree cuatro hilos; cada uno procesa un arreglo propio de calificaciones y actualiza el contador correspondiente para cada valor.
\begin{enumerate}
    \item Implemente primero la versión sin sincronización y ejecútela repetidamente con arreglos grandes.
    \item Explique por qué los resultados pueden variar aunque cada hilo procese datos distintos.
    \item Corrija el programa usando un mutex y compruebe que \texttt{aprobados + rechazados} coincide siempre con el total de datos.
\end{enumerate}

\section*{Ejercicio 4}
Lance tres hilos con una estructura que contenga un identificador, una palabra y un tiempo de espera distinto.
\begin{enumerate}
    \item Cada hilo debe construir dinámicamente una nueva cadena con la palabra invertida y almacenarla en su estructura.
    \item El \texttt{main} debe hacer \texttt{join} en el orden de los identificadores, no en el orden en que los hilos terminan.
    \item Imprima cada resultado junto con su identificador y libere las cadenas construidas.
\end{enumerate}
\end{document}
